\scanpage{50}
$U$ ouvert $\subset X$, $i,j\in I$ \quad $\overline F_i\cap U\neq\varnothing$ i.e. $F_i\cap U\neq\varnothing$
\deleted{$\unsure{\overline F_j\cap U\neq\varnothing}$}
\[
\overline F_i\subset\overline F_j
\overset{?}{\Longleftrightarrow}
\overline F_i\cap U\subset\overline F_j\cap U
\]
$\Rightarrow$ trivial

$\overset{?}{\Leftarrow}$

\deleted{\unsure{On a} $\unsure{\overline F_i\cap U\subset\overline F_j\cap U}$}

\deleted{$\unsure{F_i,F_j}$} soit $x\in\overline F_i$, au vois. de $x$ on a $\overline F_i\subset\overline F_j$,
\[
\text{i.e. }\quad
\left(\overline F_i\cap\overline F_j=\bigcup_{k\in K}\overline F_k\right)=\overline F_i
\quad\text{au vois. de }x
\]
\deleted{\unsure{alors partout}}
\[
\forall k\qquad \overline F_k\subset\overline F_i
\quad\text{égalité au voisinage de }x
\quad\Longrightarrow\quad\text{égalité partout ?}
\]
Oui, car $\overline F_k$ est \uline{\unsure{rare}} dans $\overline F_i$ !

Donc si $I_U=\{i\in I\mid\overline F_i\cap U\neq\varnothing\}\subset I$
alors la relation d'ordre de $I_U$ est celle induite par $U$.

\uline{Cor} Si $F_i\cap U\neq\varnothing$, alors $F_i$ ouvert $\Longleftrightarrow F_i\cap U$ ouvert

(i.e. $i$ maximal dans $I\Longleftrightarrow i$ maximal dans $I_U$)

\deleted{$\unsure{i}$ \unsure{clair, car si}}

$\Rightarrow$ trivial

$\Leftarrow$ soit $i$ maximal dans $I_U$, et $j\in I$, $j\geq i$, alors $j\in I_U$, donc $j=i$, OK.

\margin{\uline{III}}
\[
X=X^0\supset X^1\supset X^2\supset\cdots\supset X^n\supset\cdots
\]
\margin{\uline{filtration loc. finie par des fermés}}

$\forall i$, \deleted{$\unsure{X^i-X^1}$} $X^{\unsure{i-1|i+1}}$ rare dans $X^i$
\[
V^i=X^i-X^{\unsure{i-1|i+1}}
\qquad
\left[\text{donc }(X^{\unsure{i-1|i}}=V^i)\Longleftrightarrow X^{\unsure{i-1|i+1}}=\varnothing\right]
\]

$\forall$ ouvert $U$, l'ens. des composantes connexes de $U\cap V^i$ est loc. fini sur $U$
\uuline{et} \deleted{$\unsure{\forall}$} comp. conn. $Z_\alpha$ de $U\cap V^i$,
\inserted{$\forall j\in I$,} $\overline Z_\alpha\cap U\cap V_j$
$\longrightarrow$ une partie ouverte de $U\cap V_j$.

\margin{C'est en fait une condition locale, \unsure{sur tout ens. de} \unsure{vois. quelconques} $U_x$, $x\in X$}
\margin{51}
